% ECONOMICS_PROBLEM: {"id": "F12-381", "title": "Gravity beyond convenient CES structure", "jel_code": "F12", "source_id": "OP-0082"}
% !TeX program = xelatex
\documentclass[11pt,a4paper]{article}
\usepackage[margin=23mm]{geometry}
\usepackage{fontspec}
\setmainfont{Latin Modern Roman}
\usepackage{amsmath,amssymb,mathtools}
\usepackage{xcolor,enumitem,booktabs,longtable,array,xurl,hyperref}
\definecolor{muted}{HTML}{53636C}
\hypersetup{colorlinks=true,urlcolor=blue,linkcolor=blue}
\setlength{\parindent}{0pt}
\setlength{\parskip}{5pt}
\newcommand{\R}{\mathbb R}
\newcommand{\E}{\mathbb E}
\newcommand{\N}{\mathbb N}
\newcommand{\ind}{\mathbf 1}
\newcommand{\Var}{\operatorname{Var}}
\newcommand{\Cov}{\operatorname{Cov}}
\newcommand{\supp}{\operatorname{supp}}
\DeclareMathOperator*{\argmin}{arg\,min}
\DeclareMathOperator*{\argmax}{arg\,max}
\newcommand{\entrymeta}[3]{{\sffamily\small\color{muted}\textbf{\texttt{#1}}\quad|\quad #2\par #3\par}}
\newcommand{\Problemref}[1]{\texttt{#1}}
\newcommand{\JELref}[2]{\texttt{#2} (JEL \texttt{#1})}
\begin{document}
\section*{Gravity beyond convenient CES structure}
\textbf{Problem ID:} \texttt{F12-381}\quad \textbf{JEL:} \texttt{F12}\par
% BEGIN ECONOMICS STATEMENT
\label{problem:OP-0082}
\label{atlas:prob:T2}

\noindent\textbf{Atlas ID:} T2; \textbf{original number:} 82; \textbf{field:} International trade. \textbf{Classification:} Editorial research formulation (theoretical/empirical); not a certified unsolved theorem.

\subsection*{Definitions and assumptions}
Consider finitely many countries, products and consumer types. Type $h$ has strictly increasing preferences over a nonnegative consumption vector, with expenditure function $a_h(p,u)$ at strictly positive prices $p$ and utility $u$. Firms face specified technologies, fixed export costs and possibly variable markups. Model $\theta$ includes demand, technology, ownership, trade wedges and an equilibrium selection in which consumers maximize utility within budgets, firms maximize profits within their production sets, and all goods and factor markets clear. Let $P_O$ denote the population law of product-level prices, quantities, expenditure shares and firm characteristics under observed policies.

\subsection*{Formal research question}
For a specified policy perturbation from $\tau^0$ to $\tau^1$, write $(p^j_\theta,m^j_{h,\theta})$ for equilibrium prices and disposable expenditure, $j\in\{0,1\}$, and $v_h(p,m)$ for indirect utility. Define baseline-price equivalent variation by
\[
 \mathrm{EV}_{h,\theta}
 =a_h(p^0_\theta,v_h(p^1_\theta,m^1_{h,\theta}))-m^0_{h,\theta}.
\]
Within an explicit model class $\Theta$, characterize the range of $\mathrm{EV}_{h,\theta}$ among models that reproduce $P_O$. Derive restrictions under which bilateral flows and a limited set of elasticities identify this variation, and construct observationally equivalent models showing failure when those restrictions are removed. Determine whether product prices, household demand responses or firm markup observations restore identification.

\subsection*{Interpretation and scope}
This asks how strongly welfare conclusions depend on demand and supply structure, rather than whether a gravity regression can fit trade. An expenditure function requires well-behaved preferences and finite minimum expenditure; these are maintained assumptions to verify for each proposed specification. Income changes include profits, tariff rebates and factor returns, so holding nominal income fixed is a different counterfactual. Nonhomothetic preferences make welfare depend on the household’s initial expenditure. Theoretical nonidentification and empirical robustness are separate outputs; a finite collection of successful calibrations does not exhaust the compatible model class.

\subsection*{Source audit and status}
[S1] supplies structural gravity with multilateral resistance; [S2] supplies heterogeneous-firm margins; [S3] measures unequal gains using nonhomothetic demand. None implies universal robustness to flexible demand and endogenous markups. The question is an editorial identification and model-comparison agenda, with no exhaustive certification of its current unresolved status.

\subsection*{References}

\begin{enumerate}[label={[S\arabic*]},leftmargin=*]

\item James E. Anderson and Eric van Wincoop (2003). \emph{Gravity with Gravitas: A Solution to the Border Puzzle}. American Economic Review 93(1), 170–192. \url{https://doi.org/10.1257/000282803321455214}.
\textit{Locator and role:} Publisher or author abstract / bibliographic record; Structural gravity and multilateral resistance. Provides background or a benchmark result; does not establish that the editorial question remains unresolved in 2026.

\item Thomas Chaney (2008). \emph{Distorted Gravity: The Intensive and Extensive Margins of International Trade}. American Economic Review 98(4), 1707–1721. \url{https://doi.org/10.1257/aer.98.4.1707}.
\textit{Locator and role:} Publisher or author abstract / bibliographic record; Gravity with heterogeneous firms and extensive and intensive margins. Provides background or a benchmark result; does not establish that the editorial question remains unresolved in 2026.

\item Pablo D. Fajgelbaum and Amit K. Khandelwal (2016). \emph{Measuring the Unequal Gains from Trade}. Quarterly Journal of Economics 131(3), 1113–1180. \url{https://www.nber.org/papers/w20331}.
\textit{Locator and role:} Publisher or author abstract / bibliographic record; Nonhomothetic demand and distributional gains; NBER record identifies the 2016 published article. Provides background or a benchmark result; does not establish that the editorial question remains unresolved in 2026.

\end{enumerate}

\subsection*{Common conventions: Source mathematical conventions}

Unless an entry states otherwise, agent and firm sets are finite. State and action spaces are measurable, strategies and outcome maps are measurable, probability laws are specified, and expectations exist for the targets under discussion. Infinite-horizon discount factors lie in $(0,1)$ and discounted objectives are finite. Norms, tolerances and complexity input sizes must be specified for computational comparisons. Constants or primitives that appear locally are inputs, not empirical facts asserted by this catalogue.

\textbf{Identification.} A statistical model $m\in\mathcal M$ implies an observable law $P_m$ and target $\theta(m)$. At law $P$, its identified set is
\[
\Theta(P;\mathcal M)=\{\theta(m):m\in\mathcal M,\ P_m=P\}.
\]
Point identification holds when this set is a singleton, or equivalently when $P_{m_1}=P_{m_2}$ implies $\theta(m_1)=\theta(m_2)$ for all compatible models. Sharp bounds mean bounds attained, or approached when the boundary is not attained, by compatible models. Writing this set is a definition, not a solution: the substantive task is to establish informative restrictions and derive or estimate the set. An unrestricted-confounding model may give only trivial bounds.

\textbf{Inference.} Identification, estimability and valid uncertainty quantification are separate questions. Fix $\alpha\in(0,1)$. A confidence procedure $C_n$ over a declared law class $\mathcal P$ has asymptotic uniform coverage at level $1-\alpha$ when
\[
\liminf_{n\to\infty}\inf_{P\in\mathcal P}\Pr_P\{\theta(P)\in C_n\}\geq1-\alpha.
\]
For set-identified targets the coverage event must be defined separately, for example coverage of the full identified set. If an entry uses identification-indexed models instead of uniquely indexed laws, the same coverage convention is applied over those models. Pointwise limits do not establish uniform validity, and asymptotic validity does not guarantee finite-sample precision. Sample sizes, dependence structures and weak-identification sequences are part of each application.

\textbf{Counterfactuals.} $Y_i(a)$ is a potential outcome under a specified intervention, including its timing, implementation and versions. It is not $Y_i$ conditional on voluntarily choosing $a$. A full intervention vector permits interference; shorthand scalar treatment notation does not silently impose no interference. Assignment, selection, attrition, anticipation and measurement must be specified. A claim about an instrument's local effect is not automatically a population policy effect.

\textbf{Economic equilibrium.} A counterfactual model includes best responses, constraints, feasibility, market clearing, state transitions and expectations consistent with the stated information. When several equilibria exist, either a declared selector is an input or the target is set-valued. Comparisons across policy regimes must use the stated selection convention. Reduced-form relationships are not presumed invariant after an equilibrium-changing intervention.

\textbf{Optimization and welfare.} Optimization is over the stated feasible set. If attainment is not established, $\sup$ or $\inf$ and $\varepsilon$-optimal choices replace an asserted maximizer or minimizer. For example, with value $V=\sup_{a\in\mathcal A}W(a)<\infty$, an $\varepsilon$-optimal set is $\{a\in\mathcal A:W(a)\geq V-\varepsilon\}$ for $\varepsilon>0$. Benefits, costs, risk and health or environmental outcomes can be aggregated only using declared units and conversion weights. Interpersonal utility weights, the treatment of fiscal transfers and foreign profits, discounting, and the behavioral welfare criterion are explicit inputs. A comparative-static derivative additionally requires existence and appropriate differentiability of the selected counterfactual.

\textbf{Sources.} Local labels [S1], [S2], and so forth restart in every question. Locators identify the relevant abstract, model, section or result at the level actually checked; a bibliographic or abstract-level verification is not represented as inspection of an entire proof. Working papers are distinguished from later publications and changing versions. Source summaries are paraphrased. All URL checks and editorial formulations refer to the audit date stated above.
% END ECONOMICS STATEMENT
\end{document}
